% © 2026 Ing. David Jaroš — CC BY-NC-ND 4.0
%
% This work is licensed under a Creative Commons Attribution-NonCommercial-NoDerivatives
% 4.0 International License (CC BY-NC-ND 4.0).
%
% License History: Earlier drafts (up to v0.3) were released under CC BY 4.0.
% From v0.4 onward, all material is released under CC BY-NC-ND 4.0 to protect
% the integrity of the theoretical work during ongoing academic development.

% UBT-AI-PROVENANCE-BEGIN
% schema: ubt-ai-provenance/v1
% tier: C_working
% ai_assistance: disclosed
% human_review: risk-based
% editorial_responsibility: Ing. David Jaroš
% policy: ../../AI_PROVENANCE.md
% notice: Working material; exhaustive human review is not claimed.
% UBT-AI-PROVENANCE-END

\documentclass[11pt,a4paper]{article}
\usepackage{amsmath,amssymb,amsthm}
\usepackage{booktabs}
\usepackage{geometry}
\usepackage{hyperref}
\geometry{margin=2.5cm}

\newtheorem{proposition}{Proposition}[section]
\newtheorem{remark}[proposition]{Remark}
\newtheorem{definition}[proposition]{Definition}

\title{U(1) Generator Normalization in UBT\\
\large Target 1 — fix\_or\_reject\_U1\_coupling\_normalization}
\author{Ing.\ David Jaroš}
\date{2026-05-10}

\begin{document}
\maketitle

\begin{abstract}
We determine whether the electromagnetic U(1) generator is fixed by the
biquaternion algebra $\mathcal{B}=\mathbb{C}\otimes_{\mathbb{R}}\mathbb{H}$,
compute its trace normalization, compare it with standard conventions, and
assess whether the unit charge $q=1$ and charge quantization are derived or
assumed.  Verdict: \textbf{CONDITIONAL}.
\end{abstract}

% ----------------------------------------------------------------
\section{Identification of the U(1) generator in the UBT algebra}
% ----------------------------------------------------------------

\subsection{Biquaternion algebra and its automorphism group}

The canonical UBT algebra is
\begin{equation}
  \mathcal{B}=\mathbb{C}\otimes_{\mathbb{R}}\mathbb{H}
  \;\cong\;\mathrm{Mat}(2,\mathbb{C}).
\end{equation}
Its real automorphism group is
\begin{equation}
  \mathrm{Aut}_{\mathbb{R}}(\mathcal{B})
  \;\cong\;\bigl[\mathrm{GL}(2,\mathbb{C})\times\mathrm{GL}(2,\mathbb{C})\bigr]/\mathbb{Z}_2.
\end{equation}
Restricting to the unitary subgroup:
\begin{equation}
  \mathrm{Aut}_{\mathbb{R}}(\mathcal{B})\big|_{\mathrm{unitary}}
  \;\supset\;\mathrm{SU}(2)_L\times\mathrm{U}(1)_Y
  \quad(\text{left factor})
  \;\oplus\;\mathrm{U}(1)_{\mathrm{EM}}\quad(\text{right factor}).
\end{equation}

\subsection{The electromagnetic U(1) as the right-phase action}

The right-multiplication map
\begin{equation}
  \Theta(x)\;\longmapsto\;\Theta(x)\,e^{i\chi(x)},
  \qquad \chi(x)\in\mathbb{R},
\end{equation}
generates a U(1) subgroup of the right $\mathrm{GL}(2,\mathbb{C})$ factor.
In the $\mathrm{Mat}(2,\mathbb{C})$ representation this corresponds to
right-multiplication by the scalar phase $e^{i\chi}I_2$, i.e.\ the generator
is proportional to the $2\times 2$ identity:
\begin{equation}
  T_{\mathrm{EM}} = \frac{1}{2}I_2.
\end{equation}
The normalization factor $1/2$ is fixed by the requirement that
$\exp(2\pi i T_{\mathrm{EM}})=I_2$ (period $2\pi$ in $\chi$).

\subsection{Centre of the algebra}

The centre of $\mathcal{B}$ is $Z(\mathcal{B})=\mathbb{C}\cdot 1$, i.e.\
the complex scalars.  The electromagnetic U(1) is precisely the unitary part
of the centre:
\begin{equation}
  \mathrm{U}(1)_{\mathrm{EM}}\;=\;\bigl\{e^{i\chi}\cdot I_2\;\bigr|\;\chi\in[0,2\pi)\bigr\}
  \;\subset\; Z(\mathcal{B}).
\end{equation}

% ----------------------------------------------------------------
\section{Trace normalization computation}
% ----------------------------------------------------------------

\begin{definition}[Generator trace normalization]
For a generator $T$ in a representation $\mathcal{R}$, define
$T(\mathcal{R}) := \mathrm{Tr}_{\mathcal{R}}(T^2)$.
\end{definition}

\subsection{In the Mat(2,C) (fundamental) representation}

With $T_{\mathrm{EM}}=\tfrac{1}{2}I_2$:
\begin{equation}
  \mathrm{Tr}\!\left(T_{\mathrm{EM}}^2\right)
  =\mathrm{Tr}\!\left(\tfrac{1}{4}I_2\right)
  =\frac{2}{4}=\frac{1}{2}.
\end{equation}

\subsection{Comparison table}

\begin{center}
\begin{tabular}{@{}llll@{}}
\toprule
Convention & Generator & $\mathrm{Tr}(T^2)$ & Basis \\
\midrule
QED ($e$ absorbed into $\mathcal{A}_\mu$)
  & $q_r$ (charge ratio) & — & extrinsic \\
Standard Model hypercharge $U(1)_Y$
  & $Y/2$ with $\mathrm{Tr}(Y^2/4)=1/2$
  & $1/2$ & GUT input \\
GUT normalization $U(1)$ inside SU(5)
  & $T_{\mathrm{GUT}}=\sqrt{3/5}\,Y/2$
  & $3/10$ & rescaled \\
UBT right-phase $U(1)_{\mathrm{EM}}$
  & $\tfrac{1}{2}I_2$
  & $\mathbf{1/2}$ & algebraic \\
\bottomrule
\end{tabular}
\end{center}

\textbf{Key finding:} The UBT algebra alone yields
$\mathrm{Tr}(T_{\mathrm{EM}}^2)=1/2$, which matches the standard hypercharge
normalization.  The GUT normalization differs by a factor of $3/5$ and requires
an embedding choice that is external to $\mathcal{B}$.

% ----------------------------------------------------------------
\section{Is the unit charge $q=1$ derived or assumed?}
% ----------------------------------------------------------------

The charge operator in UBT is $Q=T_3+Y/2$ after electroweak symmetry breaking.
With the conventions above, the minimum charge is set by the condition that
$\Theta$ carries one unit of U(1) charge under the right-phase action:
\begin{equation}
  \Theta\;\longmapsto\; e^{i\chi}\Theta
  \quad\Rightarrow\quad q_\Theta = 1.
\end{equation}
This means the unit charge is the charge of the fundamental field $\Theta$.
However, the \emph{physical} unit of electric charge (the electron charge $e$)
is then related by
\begin{equation}
  e = g\sin\theta_W,
\end{equation}
which introduces the electroweak mixing angle $\theta_W$.
The Weinberg angle is not fixed by the biquaternion algebra alone: it requires
the ratio of SU(2)$_L$ and U(1)$_Y$ couplings, which are separate parameters.

\begin{remark}[CONDITIONAL status of charge unit]
The biquaternion algebra fixes the \emph{generator normalization}
$\mathrm{Tr}(T_{\mathrm{EM}}^2)=1/2$ and the charge of the fundamental field
$q_\Theta=1$.  The physical electron charge $e=g\sin\theta_W$ is derived from
$\mathcal{B}$ only if the Weinberg angle is also derived, which is an additional
open problem.
\end{remark}

% ----------------------------------------------------------------
\section{Charge quantization from topology}
% ----------------------------------------------------------------

On the compact imaginary-time circle $S^1_\psi$ of circumference $2\pi R_\psi$,
single-valuedness of $\Theta(x,\psi+2\pi R_\psi)=\Theta(x,\psi)$ under a
holonomy $\exp(iq\oint A_\psi\,d\psi)=1$ forces
\begin{equation}
  q\,\Phi_\psi = 2\pi n,\qquad n\in\mathbb{Z},
\end{equation}
where $\Phi_\psi$ is the total flux through the compact cycle.
If the minimal flux is normalized to $\Phi_\psi=2\pi$ (unit flux quantum),
then $q=n\in\mathbb{Z}$: charge quantization in integer units follows from the
winding condition.  Fractional charges (quarks, $q=1/3$) require either
a GUT-type embedding with a larger gauge group or an additional argument about
multi-valuedness in a covering space.

\begin{remark}[CONDITIONAL status of fractional charges]
Integer charge quantization is a topological consequence of the
$S^1_\psi$ compactification.  Fractional quark charges are not derived
from this topology alone.
\end{remark}

% ----------------------------------------------------------------
\section{Summary and verdict}
% ----------------------------------------------------------------

\begin{center}
\begin{tabular}{@{}lll@{}}
\toprule
Item & Status & Notes \\
\midrule
U(1) identification in $\mathcal{B}$ & DERIVED & Right-phase of $\mathrm{Mat}(2,\mathbb{C})$ \\
$\mathrm{Tr}(T_{\mathrm{EM}}^2)=1/2$ & DERIVED & Algebraic, no external input \\
QED normalization match & DERIVED & Agrees with SM hypercharge convention \\
GUT normalization match & CONDITIONAL & Requires SU(5) embedding choice \\
Unit charge $q_\Theta=1$ & DERIVED & From right-phase period \\
Electron charge $e=g\sin\theta_W$ & CONDITIONAL & Requires $\theta_W$ input \\
Integer charge quantization & CONDITIONAL & Follows from $S^1_\psi$ winding if scale fixed \\
Quark fractional charges & CONDITIONAL & Requires GUT or additional embedding \\
\bottomrule
\end{tabular}
\end{center}

\paragraph{Overall verdict.}
\textbf{CONDITIONAL.}
The U(1) generator is identified unambiguously within $\mathcal{B}$, its trace
normalization is fixed algebraically to $1/2$, and the fundamental field
$\Theta$ carries charge $q=1$ by construction.  However, the physical electron
charge $e=g\sin\theta_W$ is conditional on the Weinberg angle, and fractional
quark charges require additional embedding assumptions.  Therefore, generator
normalization is derived, but the physical charge unit remains conditional.

\end{document}
